% =============================================================================
% Бүлэг 1: Холбоотой ажлуудын судалгаа
% =============================================================================

\chapter{Холбоотой ажлуудын судалгаа}
\label{ch:related}
\pagecolor{white}

Энэ бүлэгт Text2SQL, мэдлэг хайлт, олон хэлт ба бага нөөцтэй хэлний чиглэлээр хийгдсэн судалгаа болон зах зээл дээрх
ижил төстэй бүтээгдэхүүнийг тоймлож, тэдгээрийн арга зүй, хязгаарлалтыг харьцуулан шинжилнэ. Ингэснээр энэхүү ажлын
онол-аргазүйн суурийг тодорхойлж, шийдвэрлэх судалгааны цоорхойг илрүүлнэ.

\section{Холбоотой судалгаанууд}
\label{sec:similar}

\subsection{Text2SQL даалгавар ба жишиг үнэлгээ}

Shi нар Text2SQL-ийн хөгжлийг дүрэмд суурилсан, Seq2Seq, урьдчилан сургасан хэлний загвар, LLM гэсэн үе шатаар
тодорхойлсон \cite{shi_2025_survey_llm_text2sql}. Spider \cite{yu_2018_spider} олон хүснэгттэй, урьд үзээгүй схемд
ерөнхийлөх даалгаврыг, BIRD \cite{li_2023_bird} бодит «бохир» өгөгдөл, том схем, гадаад мэдлэг (evidence) шаарддаг
нөхцөлийг тавьсан. Нэг асуултад олон зөв SQL байж болох тул текстээр бус гүйцэтгэлийн үр дүнгээр (EX) үнэлдэг; санамсаргүй
таарахаас сэргийлэхийн тулд өгөгдлийг өөрчилсөн хэд хэдэн хуулбар DB дээр шалгах test-suite аргыг санал болгосон
\cite{zhong_2020_test_suite}. TrustSQL нь хариулах боломжгүй асуултыг оруулж, зөв татгалзалтыг зөв хариулттай адил
тооцдог \cite{lee_2024_trustsql}; AMBROSIA, AmbiQT нь нэг асуулт олон зөв тайлбартай байж болохыг харуулсан
\cite{saparina_2024_ambrosia, bhaskar_2023_ambiqt}.

\subsection{Нээлттэй эхийн тусгай Text2SQL загварууд}

OmniSQL нь 2.5 сая синтетик хосоор Qwen2.5-Coder-ийг сургаж, 7B загварыг том ерөнхий загваруудтай өрсөлдөхүйц болгосон
\cite{li_2025_omnisql, hui_2024_qwen25_coder}. Arctic-Text2SQL-R1 нь OmniSQL-оос эхлүүлж, зөвхөн гүйцэтгэлд суурилсан
энгийн reward-аар (1: зөв, 0.1: ажилласан ч буруу, 0: алдаа) GRPO \cite{shao_2024_deepseekmath} хийсэн. Түүний 32B
хувилбар BIRD-ийн самбарт 71.83\%-иар тэргүүлсэн бол энэ ажилд ашигласан 7B хувилбар BIRD test дээр 68.5\% EX авсан
\cite{yao_2025_arctic_text2sql_r1}. Эдгээр загвар англи өгөгдлөөр сургагдсан тул англи бус асуултад хэр ажиллах нь
тодорхойгүй хэвээр байна.

\subsection{Мэдлэг хайлт ба жишээ сонголт}

RAG нь загварын параметрт байхгүй мэдлэгийг хайлтаар олж prompt-д нэмдэг \cite{lewis_2020_rag}. Text2SQL-д схемийн
тайлбар, хадгалагдсан утга, бизнес нэр томьёо, төстэй асуулт--SQL жишээг хайдаг. CHESS нь утгыг LSH-ээр
\cite{talaei_2024_chess}, CodeS нь BM25 ба LCS-ээр \cite{li_2024_codes} хайдаг. DAIL-SQL нь асуултын утгыг далдалж,
бүтцээр нь жишээ сонгох нь few-shot-ийн чанарыг сайжруулдгийг харуулсан \cite{gao_2024_dail_sql}. Schema linking нь
хүчтэй загварт ашиггүй ч жижиг загварт EX-ийг өсгөж болно \cite{maamari_2024_death_schema_linking}.

\subsection{Олон хувилбар, засвар ба найдвартай байдал}

Self-consistency нь хэд хэдэн хариу үүсгэж олонхоор сонгоно \cite{wang_2023_self_consistency}; Text2SQL-д нэр дэвшигч
SQL-ийг гүйцэтгээд үр дүнгээр санал хураадаг. CHASE-SQL олон генератор ба сонгогчийг хослуулах нь ганц генератороос илүү
болохыг харуулсан \cite{pourreza_2024_chase_sql}. Харин гадаад дохиогүй self-correction нь заримдаа зөв хариултыг эвддэг
\cite{huang_2024_self_correct}.

\subsection{Олон хэлт ба бага нөөцтэй хэлний Text2SQL}

MultiSpider нь Spider-ийг 6 хэл рүү орчуулж англи бус хэл дээр нарийвчлал буурдгийг \cite{dou_2023_multispider},
MultiSpider 2.0 дээр хамгийн сайн загварууд ч 4\% EX авсныг \cite{pham_2026_multispider2} харуулсан. IndicDB нь латинаар
бичсэн, англитай холилдсон асуултын нөлөөг \cite{dawar_2026_indicdb}, TUR2SQL ба BIRDTurk нь турк хэлийг
\cite{kanburoglu_2023_tur2sql, aktas_2026_birdturk}, Dialect2SQL нь араб аялгууг \cite{chafik_2025_dialect2sql} судалсан.
Translate-test стратеги (асуултыг англи руу орчуулж англи төвтэй загварт өгөх) ихэвчлэн үр дүнтэй
\cite{etxaniz_2024_self_translate, huang_2023_xlt} боловч Text2SQL-д нэр, утга гажих эрсдэлтэй \cite{dawar_2026_indicdb}.
Монгол хэлний хувьд MM-Eval нь загваруудын утга зүй, сэтгэн бодох чадвар сул \cite{zhang_2024_mm_eval}, MonCulture-Eval нь
бичиг хоорондын зөрүү байгааг \cite{minggad_2026_monculture_eval} харуулсан. Өмнөх дипломын ажил
\cite{text2sql_thesis_2026} өөр схем дээр Planner → Writer → Critic агентын урсгалыг судалсан.

\subsection{Үйлдвэрлэлийн системүүд}

Uber QueryGPT нь workspace ба тусгай агентуудтай, бүрэлдэхүүн бүрийг тусад нь үнэлдэг \cite{uber_querygpt}. LinkedIn SQL
Bot нь мэдлэгийн граф ашиглаж, асуулгыг EXPLAIN-ээр шалгадаг \cite{linkedin_sqlbot}. Эдгээр нь Text2SQL-ийг ганц загвар
биш, олон бүрэлдэхүүнтэй систем гэж үзэх хандлагыг харуулдаг.

\section{Ижил төстэй бүтээгдэхүүний судалгаа}
\label{sec:products}

Байгалийн хэлээр өгөгдлөөс асуулт асуух боломжийг өгөгдлийн платформ ба BI хэрэгслүүд бүтээгдэхүүн болгон санал болгож
байна. Тэдгээрийн мэдлэгээ хэрхэн авдаг, хаана ажилладаг, монгол хэл дэмжих эсэхийг судлав (албан ёсны баримт бичгээс,
2026 оны 10-р сарын байдлаар).

\subsection{Олон улсын арилжааны бүтээгдэхүүн}

\textbf{Snowflake Cortex Analyst.} Snowflake-ийн өгөгдлийн агуулах дээр ажиллах Text2SQL үйлчилгээ. Хүснэгт, баганын
тайлбар, синоним, баталгаажсан жишээ асуулгыг YAML хэлбэрийн \emph{semantic model}-д бичиж өгдөг; зөвхөн схем өгөхөд
ерөнхий Text2SQL алддаг асуудлыг энэ загвар нөхдөг гэж баримт бичигт тайлбарласан \cite{snowflake_cortex_analyst}.

\textbf{Databricks AI/BI Genie.} Нэг сэдвийн хүснэгтүүд, ерөнхий заавар, жишээ SQL асуулга, SQL функцээс бүрдэх «Genie
space» үүсгэж, асуултад тохирох жишээг сонгон SQL үүсгэдэг. Нэг space-д 100 хүртэл заавар (жишээ асуулга бүр нэгд
тооцогдоно) оруулж болно \cite{databricks_genie}.

\textbf{Microsoft Power BI Q\&A ба Copilot.} Q\&A нь өгөгдлийн загварын синоним, хэл зүйн харилцаа (linguistic schema)-д
тулгуурласан байгалийн хэлний асуулт хариулт бөгөөд Copilot нь синоним санал болгож тохиргоог хөнгөвчилдөг
\cite{powerbi_qna}.

\textbf{Amazon QuickSight Q.} Нэг сэдвийн dataset-үүдийг «topic» болгож, талбарын нэр, синонимыг автоматаар бэлтгээд
бизнесийн нэмэлт тодорхойлолтыг зохиогч оруулдаг \cite{aws_quicksight_q}.

\subsection{Нээлттэй эхийн бүтээгдэхүүн}

\textbf{Vanna.} DDL, баримт бичиг, жишээ SQL-ээр «сургасан» (индексжүүлсэн) RAG-д суурилж SQL үүсгэдэг MIT лицензтэй
Python сан; локал LLM-тэй ажиллах боломжтой \cite{vanna_ai}.

\textbf{Wren AI.} Бизнесийн тодорхойлолт, жишээ, хяналтыг агуулсан semantic layer (MDL)-тэй нээлттэй эхийн «GenBI» систем
\cite{wren_ai}.

\subsection{Монгол хэлний бүтээгдэхүүн}

Монгол хэлээр өгөгдлийн сангаас асуулт асуух олон нийтэд нээлттэй бүтээгдэхүүн бидний хайлтаар олдсонгүй. Дээрх
бүтээгдэхүүнүүд олон хэлт LLM ашигладаг тул монгол асуулт хүлээн авч болох ч монгол хэл дээрх нарийвчлалын талаар
мэдээлэл нийтлээгүй бөгөөд арилжааны хувилбарууд нь өгөгдлийг тухайн үүлэн платформд байршуулахыг шаарддаг.

\begin{table}[H]
\centering\scriptsize
\caption{Ижил төстэй бүтээгдэхүүний харьцуулалт}
\label{tab:products}
\begin{tabularx}{\textwidth}{p{2.5cm}p{2.0cm}X p{1.5cm}p{1.6cm}}
\toprule
Бүтээгдэхүүн & Ажиллах орчин & Мэдлэг өгөх арга & Нээлттэй эх & Монгол хэл \\
\midrule
Cortex Analyst & Snowflake үүл & Semantic model (YAML): тайлбар, синоним, жишээ асуулга & үгүй & хэмжээгүй \\
AI/BI Genie & Databricks үүл & Заавар, жишээ SQL, SQL функц & үгүй & хэмжээгүй \\
Power BI Q\&A / Copilot & Microsoft үүл & Синоним, linguistic schema & үгүй & хэмжээгүй \\
QuickSight Q & AWS үүл & Topic, синоним, бизнес тодорхойлолт & үгүй & хэмжээгүй \\
Vanna & өөрийн сервер & DDL, баримт бичиг, жишээ SQL (RAG) & тийм (MIT) & хэмжээгүй \\
Wren AI & өөрийн сервер & Semantic layer (MDL) & тийм (Apache 2.0) & хэмжээгүй \\
\textbf{Энэ ажил} & \textbf{байгууллагын дотор, нэг GPU} & \textbf{Тайлбар, утга, нэр томьёо, жишээ (RAG) + LoRA сургалт} &
  \textbf{тийм} & \textbf{кирилл, латин; хэмжсэн} \\
\bottomrule
\end{tabularx}
\end{table}

Бүтээгдэхүүнүүдийн нийтлэг санаа нь схемээс гадна бизнесийн мэдлэгийг (тайлбар, синоним, жишээ асуулга) тусад нь
бүрдүүлж өгөх явдал бөгөөд энэ нь бидний мэдлэгийн сангийн хандлагыг баталж байна. Харин тэдгээр нь (1) монгол хэл дээр
нарийвчлалаа хэмжээгүй, (2) ихэнх нь гадны үүлэн платформд өгөгдөл байршуулахыг шаарддаг, (3) хариулах боломжгүй асуултаас
татгалзах, тодруулга асуух чадвараа тоогоор үнэлээгүй.

\section{Харьцуулсан шинжилгээ}
\label{sec:comparison}

Хүснэгт~\ref{tab:related-compare}-т гол судалгаануудыг энэ ажилтай харьцуулав.

\begin{table}[H]
\centering\scriptsize
\caption{Холбоотой судалгааны харьцуулалт}
\label{tab:related-compare}
\begin{tabularx}{\textwidth}{p{2.6cm}p{1.6cm}p{1.6cm}p{1.2cm}p{1.2cm}p{1.3cm}X}
\toprule
Ажил & Хэл & Загвар & RAG & Сургалт & Татгалзах & Онцлог \\
\midrule
DAIL-SQL \cite{gao_2024_dail_sql} & англи & GPT-4 & жишээ & --- & --- & Жишээ сонголт \\
CHESS \cite{talaei_2024_chess} & англи & хаалттай & утга, схем & --- & --- & Утгын хайлт (LSH) \\
OmniSQL \cite{li_2025_omnisql} & англи & 7--32B нээлттэй & --- & SFT & --- & 2.5 сая синтетик хос \\
Arctic-R1 \cite{yao_2025_arctic_text2sql_r1} & англи & 7--32B нээлттэй & --- & GRPO & --- & Гүйцэтгэлийн reward \\
TrustSQL \cite{lee_2024_trustsql} & англи & олон & --- & --- & тийм & Хариулах боломжгүй асуулт \\
MultiSpider \cite{dou_2023_multispider} & 7 хэл & олон & --- & SFT & --- & Олон хэлний зөрүү \\
BIRDTurk \cite{aktas_2026_birdturk} & турк & олон & evidence & SFT & --- & Translate-train \\
IndicDB \cite{dawar_2026_indicdb} & энэтхэг хэлнүүд & олон & hint & --- & --- & Галиг, холимог хэл \\
\textbf{Энэ ажил} & \textbf{монгол} (кирилл, латин) & \textbf{7B нээлттэй} & \textbf{схем, утга, нэр томьёо, жишээ} &
  \textbf{LoRA} & \textbf{тийм + тодруулга} & \textbf{Орчуулга + RAG + санал, нэг GPU} \\
\bottomrule
\end{tabularx}
\end{table}

\section{Судалгааны цоорхой}
\label{sec:gap}

\begin{enumerate}
  \item \textbf{Монгол Text2SQL-ийн өгөгдөл ба туршилт байхгүй.} Монгол хэлний үнэлгээний сан, хариулах боломжгүй ба
        хоёрдмол асуулттай найдвартай байдлын хэмжилт хийгдээгүй; бүтээгдэхүүнүүд ч монгол хэл дээрх нарийвчлалаа нийтлээгүй.
  \item \textbf{Бүрэлдэхүүний нөлөө тодорхойгүй.} RAG, орчуулга, нэмэлт сургалт, санал хураалтын аль нь бага нөөцтэй хэл ба
        жижиг загварт хамгийн их нөлөөтэйг нэг схем, нэг үнэлгээн дээр харьцуулсан ажил ховор.
  \item \textbf{Орчуулгын эрсдэл.} Translate-test утгыг гажуудуулдаг гэж дурдагддаг ч Text2SQL-ийн алдааны хэдэн хувь нь
        орчуулгаас үүдэлтэйг монгол хэлэнд хэмжээгүй.
  \item \textbf{Нөөц ба нууцлал.} Ихэнх шийдэл хүчтэй GPU, хаалттай API эсвэл гадны үүл шаарддаг; нэг GPU, байгууллагын
        дотор ажиллах давтагдахуйц систем дутмаг.
\end{enumerate}

\section{Бүлгийн дүгнэлт}
\label{sec:ch2-summary}

Нээлттэй 7B загварууд англи Text2SQL-д өндөр түвшинд хүрсэн ч англи бус хэл дээр мэдэгдэхүйц буурдаг. Судалгаа ба
бүтээгдэхүүн аль аль нь домэйны мэдлэгийг (тайлбар, синоним, жишээ) тусад нь бүрдүүлэх нь чухал болохыг харуулж байгаа
боловч монгол хэл дээр хэмжсэн, байгууллагын дотор ажиллах шийдэл алга. Дараагийн бүлэгт орчуулга, мэдлэг хайлт, нэмэлт
сургалт, санал хураалтыг нэгтгэсэн, бүрэлдэхүүн бүрийг тусад нь хэмжих боломжтой архитектурыг санал болгоно.
